% Part: computability
% Chapter: computability-theory
% Section: non-comp-set

\documentclass[../../../include/open-logic-section]{subfiles}

\begin{document}

\olfileid{cmp}{thy}{ncp}
\olsection{توجد مجموعات غير قابلة للحساب}

قد ثبت أن قابلية المجموعة للحساب تستلزم قابليتها للتعداد حسابيًا. وأما
العكس فلا يصح على العموم، كما تبيّن النتيجة الآتية.

\begin{thm}\ollabel{thm:K-0}
لتكن $\OLId{latin-looped}{K}_\OLMathNumeral{0}$ المجموعة $\Setabs{\tuple{\OLId{latin-ordinary}{e}, \OLId{latin-ordinary}{x}}}{\cfind{\OLId{latin-ordinary}{e}}(\OLId{latin-ordinary}{x}) \fdefined}$.
عندئذ تكون $\OLId{latin-looped}{K}_\OLMathNumeral{0}$ قابلة للتعداد حسابيًا، ولكنها غير قابلة للحساب.
\end{thm}

\begin{proof}
أما قابلية $\OLId{latin-looped}{K}_\OLMathNumeral{0}$ للتعداد حسابيًا فتثبت بجعلها مجالًا للدالة~$\OLId{latin-ordinary}{f}$
التي تعريفها
\[
\OLId{latin-ordinary}{f}(\OLId{latin-ordinary}{z}) = \umin{\OLId{latin-ordinary}{y}}{(\len{\OLId{latin-ordinary}{z}} = \OLMathNumeral{2} \land \OLId{latin-looped}{T}((\OLId{latin-ordinary}{z})_\OLMathNumeral{0}, (\OLId{latin-ordinary}{z})_\OLMathNumeral{1}, \OLId{latin-ordinary}{y}))}.
\]
إذ متى كانت $\cfind{\OLId{latin-ordinary}{e}}(\OLId{latin-ordinary}{x})$ معرّفة عثرت $\OLId{latin-ordinary}{f}(\tuple{\OLId{latin-ordinary}{e},\OLId{latin-ordinary}{x}})$ على سجل حساب
متوقف؛ وإن لم تُعرّف $\cfind{\OLId{latin-ordinary}{e}}(\OLId{latin-ordinary}{x})$، لم تُعرّف
$\OLId{latin-ordinary}{f}(\tuple{\OLId{latin-ordinary}{e},\OLId{latin-ordinary}{x}})$. وأما إن لم يرمز $\OLId{latin-ordinary}{z}$ إلى زوج، فـ$\OLId{latin-ordinary}{f}(\OLId{latin-ordinary}{z})$
غير معرّفة كذلك.

وأما انتفاء قابلية $\OLId{latin-looped}{K}_\OLMathNumeral{0}$ للحساب فهو عين عدم قابلية مشكلة التوقف
للقرار في \olref[hlt]{thm:halting-problem}.
\end{proof}

فـ$\OLId{latin-looped}{K}_\OLMathNumeral{0}$ تجمع الأزواج $\tuple{\OLId{latin-ordinary}{e},\OLId{latin-ordinary}{x}}$ التي تحقق
$\cfind{\OLId{latin-ordinary}{e}}(\OLId{latin-ordinary}{x})\fdefined$؛ ومعنى $\tuple{\OLId{latin-ordinary}{e},\OLId{latin-ordinary}{x}}\in \OLId{latin-looped}{K}_\OLMathNumeral{0}$ أن
$\cfind{\OLId{latin-ordinary}{e}}$ تتعين عند الدخل~$\OLId{latin-ordinary}{x}$، أي يتوقف حسابها هناك. ومن ثم تسمى
«مجموعة التوقف». ونخص المجموعة
$\OLId{latin-looped}{K}=\Setabs{\OLId{latin-ordinary}{e}}{\cfind{\OLId{latin-ordinary}{e}}(\OLId{latin-ordinary}{e})\fdefined}$ باسم «مجموعة التوقف الذاتي»؛
وهي مثال معياري كثير الاستعمال لما لا يقبل القرار.

\begin{thm}\ollabel{thm:K}
مجموعة التوقف الذاتي $\OLId{latin-looped}{K}=\Setabs{\OLId{latin-ordinary}{e}}{\cfind{\OLId{latin-ordinary}{e}}(\OLId{latin-ordinary}{e})\fdefined}$
!!{c.e.} ولكنها غير قابلة للقرار.
\end{thm}

\begin{proof}
  لنفرض، طلبًا للخلف، قابلية $\OLId{latin-looped}{K}$ للقرار؛ فتكون دالتها المميزة $\Char{\OLId{latin-looped}{K}}$ قابلة
  للحساب. ونضع
  \[\OLId{latin-ordinary}{d}(\OLId{latin-ordinary}{e}) = \begin{cases}
  \OLMathNumeral{1} & \text{إذا كانت\/ $\Char{K}(e) = 0$}\\
  \fundefined & \text{في غير ذلك.}
  \end{cases}
  \]
  ثم نأخذ $\OLId{latin-ordinary}{k}$ فهرسًا لـ~$\OLId{latin-ordinary}{d}$، فيكون $\OLId{latin-ordinary}{d}\simeq\cfind{\OLId{latin-ordinary}{k}}$، ومنه
  $\OLId{latin-ordinary}{d}(\OLId{latin-ordinary}{k})\simeq\cfind{\OLId{latin-ordinary}{k}}(\OLId{latin-ordinary}{k})$. وهذا يناقض حقيقة أن
  $\OLId{latin-ordinary}{d}(\OLId{latin-ordinary}{k})\fdefined$ إذا وفقط إذا كانت $\cfind{\OLId{latin-ordinary}{k}}(\OLId{latin-ordinary}{k})\fundefined$، كما
  يقتضي تعريف~$\OLId{latin-ordinary}{d}$.

  وأما $\OLId{latin-looped}{K}$ فهي مجال $\OLId{latin-ordinary}{f}(\OLId{latin-ordinary}{x})=\umin{\OLId{latin-ordinary}{y}}{\OLId{latin-looped}{T}(\OLId{latin-ordinary}{x},\OLId{latin-ordinary}{x},\OLId{latin-ordinary}{y})}$، ومن ثم فهي !!{c.e.}.
\end{proof}

\end{document}
